% Hafta 9 -- alıştırma gövdesi (ders notu ve ayrı alıştırma kağıdı tarafından \input edilir)
\section{Çözümlü Alıştırmalar}
Her soruyu önce kendiniz deneyin, sonra çözümü okuyun.

\begin{alistirma}
\((-2)^3+(-1)^{2023}-(-4)^2+3^0\) \ işleminin sonucu kaçtır?
\end{alistirma}
\begin{cozum}
\((-2)^3=-8\) (tek kuvvet), \ \((-1)^{2023}=-1\) (\(2023\) tek), \ \((-4)^2=16\) (çift kuvvet), \ \(3^0=1\).\\
Sonuç: \(-8+(-1)-16+1=-24\).
\end{cozum}

\begin{alistirma}
\(\left(\dfrac23\right)^{-2}+\left(\dfrac12\right)^{-3}\) \ işleminin sonucu kaçtır?
\end{alistirma}
\begin{cozum}
Negatif üste taban ters çevrilir: \(\left(\dfrac23\right)^{-2}=\left(\dfrac32\right)^{2}=\dfrac94\), \quad \(\left(\dfrac12\right)^{-3}=2^3=8\).
\[\frac94+8=\frac94+\frac{32}{4}=\frac{41}{4}\]
\end{cozum}

\begin{alistirma}
\((0{,}2)^{-1}+(0{,}5)^{-2}\) \ işleminin sonucu kaçtır?
\end{alistirma}
\begin{cozum}
\(0{,}2=\frac15\), \ \(0{,}5=\frac12\) \(\Rightarrow\) \(\left(\frac15\right)^{-1}+\left(\frac12\right)^{-2}=5+2^2=5+4=9\).
\end{cozum}

\begin{alistirma}
\((2^3)^4\cdot(2^{-2})^5\cdot 8^2\) \ işleminin sonucu kaçtır?
\end{alistirma}
\begin{cozum}
Üssün üssü: \((2^3)^4=2^{12}\), \ \((2^{-2})^5=2^{-10}\), \ \(8^2=(2^3)^2=2^6\).\\
\(2^{12}\cdot 2^{-10}\cdot 2^{6}=2^{12-10+6}=2^{8}=256\).
\end{cozum}

\begin{alistirma}
\(\dfrac{3^{11}+3^{10}}{3^{10}-3^{9}}\) \ işleminin sonucu kaçtır?
\end{alistirma}
\begin{cozum}
Pay kısmında en küçük kuvvet \(3^{10}\), paydada \(3^{9}\) ortak çarpana alınır:
\[\frac{3^{10}(3+1)}{3^{9}(3-1)}=\frac{4\cdot 3^{10}}{2\cdot 3^{9}}=2\cdot 3=6\]
\end{cozum}

\begin{alistirma}
\(8^{x+1}=16^{x-1}\) \ denklemini çözünüz.
\end{alistirma}
\begin{cozum}
İki taraf \(2\) tabanında yazılır: \(8^{x+1}=2^{3(x+1)}=2^{3x+3}\), \ \(16^{x-1}=2^{4(x-1)}=2^{4x-4}\).
\[3x+3=4x-4\ \Rightarrow\ x=7\qquad(\text{Kontrol: } 8^8=2^{24}=16^6)\]
\end{cozum}

\begin{alistirma}
\(2^{x+1}+2^{x}+2^{x-1}=56\) \ denklemini çözünüz.
\end{alistirma}
\begin{cozum}
En küçük kuvvet \(2^{x-1}\) parantezine alınır: \(2^{x+1}=4\cdot 2^{x-1}\), \ \(2^{x}=2\cdot 2^{x-1}\).
\[2^{x-1}(4+2+1)=7\cdot 2^{x-1}=56\ \Rightarrow\ 2^{x-1}=8\ \Rightarrow\ x=4\qquad(\text{Kontrol: } 32+16+8=56)\]
\end{cozum}

\begin{alistirma}
\(a=2^{60}\), \ \(b=3^{40}\), \ \(c=5^{20}\) \ sayılarını küçükten büyüğe sıralayınız.
\end{alistirma}
\begin{cozum}
Üslerin EBOB'u \(20\): \ \(a=(2^3)^{20}=8^{20}\), \ \(b=(3^2)^{20}=9^{20}\), \ \(c=5^{20}\).\\
Üsler eşit olduğundan tabanlar karşılaştırılır: \(5<8<9\) \ \(\Rightarrow\) \ \(c<a<b\).
\end{cozum}

\kendinizoncesi
\section{Alıştırmalar (Kendiniz Çözün)}

\begin{enumerate}[label=\textbf{\arabic*.},leftmargin=*,itemsep=6pt]
\item \(-5^2+(-5)^2-(-2)^4\) \ işleminin sonucu kaçtır?
\item \(2^{-3}+3^{-1}-4^{-1}\) \ işleminin sonucu kaçtır?
\item \((0{,}1)^{-2}+(0{,}5)^{-1}\) \ işleminin sonucu kaçtır?
\item \(2^x\cdot 2^5\cdot 2^{-3}\) \ ifadesini tek bir üslü sayı olarak yazınız.
\item \(\dfrac{3^{12}\cdot 3^{-4}}{3^5}\) \ işleminin sonucu kaçtır?
\item \(9^3\cdot 27^{-2}\cdot 81\) \ işleminin sonucu kaçtır?
\item \(14\,500\,000\) \ ve \ \(0{,}000025\) \ sayılarını bilimsel gösterimle yazınız.
\item \((6\cdot 10^{-3})\cdot(5\cdot 10^{8})\) \ işleminin sonucunu bilimsel gösterimle yazınız.
\item \(2^x=a\) olduğuna göre \(8^{x+1}\) ifadesini \(a\) türünden yazınız.
\item \(3^{2x-1}=243\) \ denklemini çözünüz.
\item \(\left(\dfrac13\right)^{x-1}=9^{x+2}\) \ denklemini çözünüz.
\item \(3^{x+2}-3^{x}=72\) \ denklemini çözünüz.
\item \(a=4^{12}\), \ \(b=3^{18}\), \ \(c=2^{30}\) \ sayılarını küçükten büyüğe sıralayınız.
\item \(x\) tam sayı olmak üzere \((x-1)^{2x}=(x-1)^{x+4}\) \ denklemini sağlayan kaç farklı \(x\) değeri vardır?
\end{enumerate}

\subsection*{Cevap Anahtarı}
{\small
\begin{multicols}{2}
\begin{enumerate}[label=\textbf{\arabic*.},leftmargin=*,itemsep=1pt]
\item \(-25+25-16=-16\).
\item \(\frac18+\frac13-\frac14=\frac{3+8-6}{24}=\frac{5}{24}\).
\item \(10^2+2=102\).
\item \(2^{x+5-3}=2^{x+2}\).
\item \(3^{12-4-5}=3^3=27\).
\item \(3^{6}\cdot 3^{-6}\cdot 3^{4}=3^4=81\).
\item \(1{,}45\cdot 10^{7}\); \ \(2{,}5\cdot 10^{-5}\).
\item \(30\cdot 10^{5}=3\cdot 10^{6}\).
\item \(8\cdot 8^x=8\cdot(2^x)^3=8a^3\).
\item \(3^{2x-1}=3^5\Rightarrow 2x-1=5\Rightarrow x=3\).
\item \(3^{-x+1}=3^{2x+4}\Rightarrow -x+1=2x+4\Rightarrow x=-1\).
\item \(3^x(9-1)=72\Rightarrow 3^x=9\Rightarrow x=2\).
\item \(a=2^{24}\); EBOB \(6\): \(16^6<27^6<32^6\Rightarrow a<b<c\).
\item \(x=4\) (üsler eşit), \(x=2\) (taban \(1\)), \(x=1\) (\(0^2=0^5\)), \(x=0\) (\((-1)^0=(-1)^4\)): \(4\) değer.
\end{enumerate}
\end{multicols}}
